\section{ANTLR}
Náš kompilátor používa generátor \textit{parserov} \textbf{AN}other \textbf{T}ool for
\textbf{L}anguage \textbf{R}ecognition. V tejto časti si priblížime, ako funguje \textit{ANTLR} a
ako prebieha preklad z Datalogu do relačnej algebry.

\subsection{Rýchlokurz formálnych jazykov}

\begin{definition}
    Pod \textbf{abecedou} rozumieme ľubovoľnú \textit{neprázdnu konečnú} množinu symbolov (znakov).
\end{definition}

\begin{definition}
    \textbf{Slovo} je \textit{konečná} postupnosť znakov (z nejakej abecedy $\Sigma$). Prázdne slovo
    sa označuje $\varepsilon$.
\end{definition}

\begin{definition}
    \textbf{Jazyk} nad abecedou $\Sigma$ je ľubovoľná množina slov nad $\Sigma$.
\end{definition}

\begin{definition}
    \textbf{Zreťazením slov} $w_1 = a_1\dots a_n$, $w_2 = b_1\dots b_m$,\\
    $a_1, \dots, a_n, b_1, \dots, b_m \in \Sigma$ dostaneme slovo $w_1 \cdot w_2 = w_1w_2 = a_1 \dots
    a_n b_1 \dots b_m$.
\end{definition}

\begin{definition}
    \textbf{Mocnina slova} je definovaná ako $w^0 = \varepsilon, w^{n+1} = w^n w$, $\forall n \in
    \mathbb{N}$.
\end{definition}

\begin{note}
    Abecedu $\Sigma$ vieme chápať aj ako množinu všetkých jednoznakových slov nad $\Sigma$.
\end{note}

\begin{definition}
    \textbf{Zreťazenie jazykov} $L_1$ (nad abecedou $\Sigma_1$) a $L_2$ (nad abecedou $\Sigma_2$) je
    jazyk $L_1 \cdot L_2 = L_1 L_2 = \{w_1 w_2 \, | \, w_1 \in L_1, w_2 \in L_2\}$ nad abecedou
    $\Sigma_1 \cup \Sigma_2$.
\end{definition}

\begin{definition}
    \textbf{Mocnina jazyka} $L$ je definovaná ako $L^0 = \{\varepsilon\}, L^{n+1} = L^nL$, $\forall
    n \in \mathbb{N}$.
    
    $L^+ = \bigcup\limits_{n > 0} L^n$
    
    $L^* = \bigcup\limits_{n \geq 0} L^n$
\end{definition}

\begin{definition}
    \textbf{Regulárna gramatika} je štvorica $G = (N, T, P, \sigma)$, kde $N$ je abeceda
    neterminálov, $T$ je abeceda terminálov, $P \subseteq_{kon} N \times (T^* \cup T^*N)$ je
    \textit{konečná} množina prepisovacích pravidiel a $\sigma \in N$ je počiatočný neterminál.
    Prepisovacie pravidlá teda prepisujú jeden neterminál na ľubovoľné slovo z terminálov, za ktorým
    môže ale nemusí, nasledovať jeden neterminál.
\end{definition}

\begin{definition}
    \textbf{Bezkontextová gramatika} je štvorica $G = (N, T, P, \sigma)$, kde $N$ je abeceda
    neterminálov, $T$ je abeceda terminálov, $P \subseteq_{kon} N \times (N \cup T)^*$ je
    \textit{konečná} množina prepisovacích pravidiel a $\sigma \in N$ je počiatočný neterminál.
    Prepisovacie pravidlá teda prepisujú jeden neterminál na ľubovoľné slovo z terminálov a
    neterminálov.
\end{definition}

\begin{note}
    Pravidlá z $P$ budeme zapisovať v tvare $n \rightarrow w \in P$ namiesto $(n, w) \in P$.
    Pravidlá s rovnakou ľavou stranou $n \rightarrow w_1$, $n \rightarrow w_2$, \dots, $n\rightarrow
    w_i$ budeme zapisovať zjednodušene ako $n \rightarrow w_1 \, | \, w_2 \, | \, \dots \, | \,
    w_i$.
\end{note}

\begin{definition}
    \textbf{Krok odvodenia} v gramatike $G$ je binárna relácia $\Rightarrow_G$ nad $(N \cup T)^*$.
    Slová $x$, $y$ sú v relácii $\Rightarrow_G$ práve vtedy, keď $x = w_1 n w_2$, $y = w_1 s w_2$ a
    $n \rightarrow s$ je pravidlo v $P$.
\end{definition}

\begin{definition}
    \textbf{Jazyk generovaný gramatikou} $G$ je množina 
    $$L(G) = \{w \in T^* \, | \, \sigma \Rightarrow_G^*w\}\text{.}$$ \textbf{Regulárne jazyky},
    resp. \textbf{bezkontextové jazyky} sú jazyky, ktoré sú generované \textit{regulárnou}, resp.
    \textit{bezkontextovou} gramatikou.
\end{definition}

\begin{definition}
    \textbf{Odvodenie} slova $w$ (v gramatike $G$) $\sigma \Rightarrow w_1 \Rightarrow w_2
    \Rightarrow \cdots \Rightarrow w$ je postupnosť krokov odvodení. Pod \textbf{ľavým}, resp.
    \textbf{pravým krajným odvodením} rozumieme odvodenie, v ktorom sa v každom kroku odvodenia, v
    ktorom máme viacero neterminálov, vždy prepíše \textit{najľavejší},
    resp. \textit{najpravejší} neterminál.
\end{definition}

\subsection{Syntaktická analýza}
Na to, aby sme vedeli preložiť Datalog do relačnej algebry, potrebujeme vedieť identifikovať a
preložiť jeho jednotlivé jazykové prvky. Tie prvky nájde tzv. \textit{syntaktická analýza}. Ukážeme
si, ako funguje na jednoduchom príklade, v ktorom budeme prekladať aritmetické výrazy z infixovej
notácie do postfixovej notácie.

Zoberme si jednoduchú gramatiku $G$ generujúcu jazyk aritmetických výrazov v infixovej notácii:
\begin{align*}
    G &= (N, T, P, \sigma)\text{, kde}\\
    N &= \{\sigma, \alpha\}\text{,}\\
    T &= \{\texttt{+}, \cdot, \texttt{(}, \texttt{)}, 0, 1, 2, 3, 4, 5, 6, 7, 8, 9\}\text{,}\\
    P &= \{\sigma \rightarrow \texttt{(}\sigma\texttt{)} \, | \, \sigma \texttt{+} \sigma \, | \, \sigma \cdot \sigma
                      \, | \, 1 \alpha \, | \, 2 \alpha \, | \, 3 \alpha
                      \, | \, 4 \alpha \, | \, 5 \alpha \, | \, 6 \alpha
                      \, | \, 7 \alpha \, | \, 8 \alpha \, | \, 9 \alpha
                      \, | \, 0,\\
      &\quad \quad \alpha \rightarrow 0 \alpha \, | \, 1 \alpha \, | \, 2 \alpha
                      \, | \, 3 \alpha \, | \, 4 \alpha \, | \, 5 \alpha
                      \, | \, 6 \alpha \, | \, 7 \alpha \, | \, 8 \alpha \, | \, 9 \alpha \, | \, \varepsilon
                      \}\text{.}
\end{align*}

Zoberme si slovo z tohto jazyka, napríklad $(4 + 5)\cdot 123 \cdot 0$, a k nemu jedno
prislúchajúce odvodenie (odvodenie zakreslíme pomocou stromu).

\begin{center}
    \[
    \Tree
    [.$\sigma$ 
        [.$\sigma$ 
            $\texttt{(}$
            [.$\sigma$
                [.$\sigma$
                    $4$
                    [.$\alpha$
                        $\varepsilon$
                    ]
                ]
                $\texttt{+}$
                [.$\sigma$
                    $5$
                    [.$\alpha$
                        $\varepsilon$
                    ]
                ]
            ]
            $\texttt{)}$
        ] 
        $\cdot$ 
        [.$\sigma$ 
            [.$\sigma$
                $1$
                [.$\alpha$
                    $2$
                    [.$\alpha$
                        $3$
                        [.$\alpha$
                            $\varepsilon$
                        ]
                    ]
                ]
            ]
            $\cdot$
            [.$\sigma$
                $0$
            ]
        ]
    ]
    \]
    \textit{Pravé krajné odvodenie}
\end{center}

\begin{center}
    \[
    \Tree
    [.$\sigma$ 
        [.$\sigma$
            [.$\sigma$ 
                $\texttt{(}$
                [.$\sigma$
                    [.$\sigma$
                        $4$
                        [.$\alpha$
                            $\varepsilon$
                        ]
                    ]
                    $\texttt{+}$
                    [.$\sigma$
                        $5$
                        [.$\alpha$
                            $\varepsilon$
                        ]
                    ]
                ]
                $\texttt{)}$
            ] 
            $\cdot$
            [.$\sigma$
                $1$
                [.$\alpha$
                    $2$
                    [.$\alpha$
                        $3$
                        [.$\alpha$
                            $\varepsilon$
                        ]
                    ]
                ]
            ]
        ]
        $\cdot$
        [.$\sigma$
            $0$
        ]
    ]
    \]
    \textit{Ľavé krajné odvodenie}
\end{center}

Môžeme si všimnúť, že veľkosť stromu odvodenia rýchlo narastá, hlavne kvôli odvodzovaniu čísiel.
Preto sa v praxi často používa tzv. \textit{lexer} (\textit{tokenizer}). Lexer rozdelí vstupný
reťazec (slovo) na tzv. \textit{tokeny}, ktoré sú zjednodušenou reprezentáciou častí vstupného
reťazca. Tieto tokeny sú popísané pomocou regulárnych jazykov (väčšinou pomocou regulárnych výrazov,
ktoré sú ekvivalentné s regulárnymi gramatikami \cite{rovan2013regex}). Následne sa odvodenie robí
nad tokenami, čím sa zmenší veľkosť stromu odvodenia a zjednoduší sa nám odvodenie. V našom prípade
definujeme jeden token reprezentujúci celé čísla
$$\texttt{INT} = \{0\} \cup \{1, 2, 3, 4, 5, 6, 7, 8, 9\} \cdot \{1, 2, 3, 4, 5, 6, 7, 8, 9, 0\}^* \approx \mathbb{N}_0\text{.}$$

Lexer pri čítaní vstupného reťazca postupne rozpoznáva jednotlivé tokeny podľa definovaných
regulárnych výrazov.  Napríklad pri spracovaní reťazca $(4 + 5)\cdot 123 \cdot 0$ by lexer generoval
nasledujúci \textit{stream tokenov}:
$$\texttt{LPAREN}(\texttt{(}),\, \texttt{INT}(4),\, \texttt{PLUS}(\texttt{+}),\, \texttt{INT}(5),\,
\texttt{RPAREN}(\texttt{)}),\, \texttt{MULT}(\cdot),\, \texttt{INT}(123),\, \texttt{MULT}(\cdot),\,
\texttt{INT}(0)$$

Každý token obsahuje svoj typ (\texttt{INT}, \texttt{PLUS}, \texttt{LPAREN}, \dots), prípadne svoju
hodnotu.  Táto sekvencia tokenov sa potom posiela parseru, ktorý na nej vykonáva syntaktickú
analýzu.

Potom sa gramatika $G$ zjednoduší na
\begin{align*}
    G' &= (N', T', P', \beta)\text{, kde}\\
    N' &= \{\beta\}\text{,}\\
    T' &= \{\texttt{+}, \cdot, \texttt{(}, \texttt{)}, \mathtt{INT}\}\text{,}\\
    P' &= \{\beta \rightarrow \texttt{(}\beta\texttt{)} \, | \, \beta \texttt{+} \beta \, | \, \beta \cdot \beta
                      \, | \, \texttt{INT}\}\text{.}
\end{align*}



Ľavé krajné odvodenie slova $(4 + 5)\cdot 123 \cdot 0$ v gramatike $G'$ by vyzeralo nasledovne:
\begin{center}
    \[
    \Tree
    [.$\beta$ 
        [.$\beta$
            [.$\beta$ 
                $\texttt{(}$
                [.$\beta$
                    [.$\beta$
                        \qroof{$4$}.$\texttt{INT}$ 
                    ]
                    $\texttt{+}$
                    [.$\beta$
                        \qroof{$5$}.$\texttt{INT}$ 
                    ]
                ]
                $\texttt{)}$
            ] 
            $\cdot$
            [.$\beta$
                \qroof{$123$}.$\texttt{INT}$ 
            ]
        ]
        $\cdot$
        [.$\beta$
            \qroof{$0$}.$\texttt{INT}$ 
        ]
    ]
    \]
    \textit{Ľavé krajné odvodenie v gramatike $G'$}
\end{center}


Na takýto strom odvodenia môžeme aplikovať algoritmus, ktorý nám z infixovej notácie
vygeneruje postfixovú notáciu. Algoritmus funguje nasledovne. Vrcholy stromu (odvodenia)
sa transformujú podľa nasledujúcich pravidiel:

\begin{center}
    \begin{tabular}{@{}c c c@{}}
        \Tree [.$\beta$ $\texttt{(}$ [.$\beta_1$ ] $\texttt{)}$ ]
        & \raisebox{-4ex}{$\rightsquigarrow$}
        & \raisebox{-4ex}{\Tree [.$\beta_1$ ]} \\
        \\
        \Tree [.$\beta$ [.$\beta_1$ ] $\texttt{+}$ [.$\beta_2$ ] ]
        & \raisebox{-4ex}{$\rightsquigarrow$}
        & \Tree [.$\beta$ [.$\beta_1$ ] [.$\beta_2$ ] $\texttt{+}$ ] \\
        \\
        \Tree [.$\beta$ [.$\beta_1$ ] $\cdot$ [.$\beta_2$ ] ]
        & \raisebox{-4ex}{$\rightsquigarrow$}
        & \Tree [.$\beta$ [.$\beta_1$ ] [.$\beta_2$ ] $\cdot$ ] \\
        \\
        \Tree [.$\texttt{INT}$ ]
        & $\rightsquigarrow$
        & \Tree [.$\texttt{INT}$ ]
    \end{tabular}
\end{center}

Po aplikovaní týchto pravidiel na strom odvodenia dostaneme strom reprezentujúci postfixovú notáciu,
z ktorého vieme jednoducho vygenerovať výsledný výraz.

\begin{center}
    \[
    \Tree
    [.$\beta$ 
        [.$\beta$
            [.$\beta$ 
                [.$\beta$
                    \qroof{$4$}.$\texttt{INT}$ 
                ]
                [.$\beta$
                \qroof{$5$}.$\texttt{INT}$ 
                ]
                $\texttt{+}$
            ] 
            [.$\beta$
            \qroof{$123$}.$\texttt{INT}$ 
            ]
            $\cdot$
        ]
        [.$\beta$
        \qroof{$0$}.$\texttt{INT}$ 
        ]
        $\cdot$
    ]
    \]
    \textit{Odvodenie slova $4 \, 5 + 123 \cdot 0 \,\cdot$}
\end{center}

Problém nastáva, ako zo vstupného slova vygenerovať strom odvodenia. Na tento problém existuje
viacero algoritmov, ktoré sa líšia v tom, akú množinu gramatík dokážu spracovať. Jeden z týchto
algoritmov je tzv. LL($k$) algoritmus, ktorý číta vstup zľava doprava (\textbf{L}eft to right) a
vytvára ľavé krajné odvodenie (\textbf{L}eftmost derivation) a pozerá sa na $k$ nasledujúcich
tokenov. 

ANTLR je generátor parserov, ktorý na základe gramatiky, ktorú mu zadáme, vygeneruje parser (a
lexer), ktorý dokáže pomocou LL($k$) algoritmu zistiť, či dané slovo patrí do jazyka generovaného
touto gramatikou. V prípade, že existuje, vygeneruje k nemu prislúchajúci strom ľavého krajného
odvodenia.